\documentclass[10pt]{article}

\usepackage[margin=1in]{geometry}
\usepackage{amsmath,amssymb,amsfonts,amsthm,mathtools,bm}
\usepackage{booktabs}
\usepackage{graphicx}
\usepackage{array}
\usepackage{enumitem}
\usepackage{xcolor}
\usepackage{microtype}
\usepackage{hyperref}
\usepackage[backend=biber,style=authoryear,maxcitenames=2,maxbibnames=99,natbib=true]{biblatex}
\addbibresource{interactive_planning_revised.bib}

\hypersetup{
    colorlinks=true,
    linkcolor=blue,
    citecolor=blue,
    urlcolor=blue
}

\newcommand{\framework}{COWP}
\newcommand{\E}{\mathbb{E}}
\newcommand{\Prob}{\mathbb{P}}
\newcommand{\ind}{\mathbf{1}}
\newcommand{\cA}{\mathcal{A}}
\newcommand{\cC}{\mathcal{C}}
\newcommand{\cD}{\mathcal{D}}
\newcommand{\cG}{\mathcal{G}}
\newcommand{\cL}{\mathcal{L}}
\newcommand{\cM}{\mathcal{M}}
\newcommand{\cN}{\mathcal{N}}
\newcommand{\cP}{\mathcal{P}}
\newcommand{\cQ}{\mathcal{Q}}
\newcommand{\cS}{\mathcal{S}}
\newcommand{\cU}{\mathcal{U}}
\newcommand{\cZ}{\mathcal{Z}}
\newcommand{\R}{\mathbb{R}}

\newtheorem{proposition}{Proposition}
\setlength{\emergencystretch}{3em}

\title{Not All Collision-Free Plans Are Safe: Non-Coercive Feasibility for Interactive Autonomous Driving}
\author{Anonymous Authors}
\date{}

\begin{document}
\maketitle

\begin{abstract} Autonomous driving planners are typically evaluated by whether the ego vehicle remains collision-free, comfortable, and efficient. However, in interactive traffic, an ego plan can appear safe only because another road user is forced to absorb the conflict through hard braking, abrupt yielding, or surrendering a legitimate gap. We identify this overlooked failure mode as \emph{false-safe planning}: the ego rollout is collision-free, but its safety is contingent on another agent taking a high-burden ceding response. 

We propose \textbf{COWP}, a coercion-witness planning framework that certifies whether an ego candidate is safe without shifting the conflict burden to others. Instead of treating social compliance as an auxiliary cost, COWP introduces \emph{non-coercive feasibility}: a plan is feasible only if each protected-priority road user retains a sufficiently large set of low-burden safe responses; an all-critical audit is reported separately as a stricter diagnostic. To operationalize this notion, COWP constructs counterfactual natural alternatives for surrounding agents, evaluates ego-conditioned safe response sets, and rejects collision-free candidates when a coercion witness shows that safety depends on high-burden ceding. We define an evidence-gated evaluation protocol on interaction-heavy Waymax scenarios that separately measures conventional closed-loop safety, ego progress, root-wise option preservation, and burden transfer. The current manuscript reports claims only after the corresponding natural-basis, transport, selector, and reactive-agent validation gates pass. \end{abstract} \section{Introduction} \label{sec:introduction} A central goal of autonomous driving planning is to produce trajectories that are safe, comfortable, and efficient. Existing planners have made significant progress by modeling interactions between the ego vehicle and surrounding road users, either through joint prediction and planning, game-theoretic reasoning, graph-based interaction modeling, or closed-loop simulation \cite{gameformer, jia2022hdgt, wang2024socialformer, renz2022plant, dauner2023parting}. Yet most planning objectives and benchmarks still judge a candidate primarily from the ego vehicle's perspective: whether the ego remains collision-free, follows the route, avoids hard braking, and reaches its goal. 

This ego-centric view misses an important failure mode in interactive traffic. Consider an ego vehicle that cuts into a target lane and remains collision-free only because the rear vehicle in that lane brakes hard and gives up its gap. Standard collision metrics may score the ego plan as safe. Even interaction-aware prediction may regard the outcome as likely if the other vehicle is predicted to yield. However, such a plan is not truly safe in the sense of cooperative road use: the ego vehicle has not resolved the conflict, but has forced another road user to absorb it. We call this mode \emph{false-safe planning}. The plan appears safe in rollout, but its safety is contingent on another agent's high-burden ceding response. 

This failure mode is related to, but distinct from, existing notions of social planning and formal safety. Courteous or socially aware planning penalizes the cost imposed on other agents, often through a soft objective or a social value orientation term \cite{sun2018courteous, schwarting2019social, crosato2022svo}. Formal safety models such as RSS specify mathematically grounded safety distances and proper responses \cite{shalev2017rss}. Causal responsibility metrics such as feasible action-space reduction study how one agent can restrict another agent's available actions \cite{george2023fear}. These ideas provide useful ingredients, but they do not directly answer the planning-time question we study: \emph{should a collision-free ego candidate be rejected because it is safe only by forcing another road user to take a high-burden ceding response?} 

We answer this question by introducing \emph{non-coercive feasibility}. A candidate ego trajectory is non-coercively feasible if every protected-priority road user retains at least one low-burden safe response and, more importantly, if the ego does not collapse that road user's natural action space into a narrow or high-burden escape maneuver. We retain an all-critical version only as a stronger diagnostic, because an ego-priority relation should not create the same hard veto as an agent-priority or explicitly negotiated relation. This changes the role of social interaction in planning. Instead of adding another courtesy-like soft penalty, we treat safety-by-coercion as a feasibility defect. A plan that relies on another vehicle's hard braking, abrupt yielding, priority abandonment, or gap surrender is not merely less polite; it is false-safe. 

\begin{proposition}[A finite soft burden penalty is not a feasibility guarantee]
Let $k_c$ be a collision-free but coercive candidate and $k_n$ a non-coercive candidate, with ego costs $J_e(k)$ and transferred burdens $B(k)$. For any finite $\lambda>0$, minimizing $J_e(k)+\lambda B(k)$ can select $k_c$ whenever
\begin{equation}
    J_e(k_n)-J_e(k_c)>\lambda\bigl(B(k_c)-B(k_n)\bigr).
\end{equation}
By contrast, a hard non-coercive feasibility constraint rejects $k_c$ independently of its ego-utility advantage.
\end{proposition}
\noindent This elementary separation explains why COWP treats coercion as a feasibility defect rather than relying only on objective reweighting.

To instantiate this idea, we propose \textbf{COWP} (\textbf{Coercion-Witness Planning}). COWP first constructs a burden-oriented interaction graph whose edges encode not only who interacts with whom, but who would have to absorb a potential conflict. The graph proposes critical ego-agent relations and candidate coercion witnesses. For each ego candidate, COWP then generates counterfactual natural alternatives for the affected agents, including observational futures, ego-neutral futures, and priority-preserving rule-constrained alternatives. These alternatives describe how an agent could reasonably behave without being pressured by the ego candidate. COWP next constructs ego-conditioned safe response sets and evaluates whether the remaining safe responses are low-burden. If a natural low-burden behavior becomes unsafe under the ego candidate and the only safe response requires excessive burden, COWP confirms a coercion witness and rejects the candidate. 

This design has three advantages. First, it separates \emph{prediction} from \emph{certification}: even if a model predicts that another vehicle will yield, COWP asks whether the ego should rely on that yield. Second, it provides an interpretable witness rather than only a scalar score, identifying the burdened agent, conflict region, required ceding behavior, and burden mechanism. Third, it is compatible with existing candidate generators and closed-loop simulators. COWP can therefore be used as a certification layer on top of rule-based, learning-based, or hybrid planners.

We make the following contributions: \begin{itemize} 
\item We identify \emph{false-safe planning}, a collision-free but coercive failure mode in which ego safety depends on another road user's high-burden ceding response. 
\item We formulate \emph{non-coercive feasibility}, a planning-time criterion requiring protected-priority agents to retain sufficiently rich low-burden safe response sets, rather than being forced into narrow or costly evasive behavior. 
\item We propose \textbf{COWP}, a coercion-witness planning framework that combines burden-oriented interaction graphs, counterfactual natural alternatives, ego-conditioned safe response sets, and witness-based candidate rejection. \end{itemize}


\section{Related Work}

\subsection{Ego-conditioned prediction and prediction--planning coupling}
Ego-conditioned prediction methods condition the forecast of surrounding agents on a proposed ego plan. PiP predicts agent trajectories given a planned ego trajectory \parencite{song2020pip}; DTPP jointly trains ego-conditioned prediction and context-aware cost evaluation for tree policy planning \parencite{huang2024dtpp}; PPAD iterates between prediction and planning interactions \parencite{chen2024ppad}; and HYPE combines learned ego proposals, ego-conditioned occupancy prediction, and search-based refinement \parencite{yu2025hype}. These methods are strong baselines because they explicitly condition futures on ego behavior. Their limitation for our problem is that they generally evaluate candidates one by one. \framework{} instead learns a shared same-root response surface, which lets the planner compare a candidate with neighboring interventions and ask whether its safety is robust or only obtained through another agent's high-burden response.

\subsection{Interactive reasoning, game models, and behavior latents}
Interactive forecasting methods model how agents influence each other. M2I decomposes prediction into influencer and reactor roles \parencite{sun2022m2i}, while GameFormer uses hierarchical game-theoretic Transformer decoding for iterative interaction prediction and planning \parencite{huang2023gameformer}. Driver-style, social, and behavior-latent approaches encode intent, aggressiveness, or latent world state \parencite{hu2023formulating,chen2025categorical,deng2025social,hao2026styledrive,zheng2025world4drive}. These representations are useful for describing what an agent may do in a scene. Our focus is different: for the same root scene, we model how the agent's response distribution changes as the ego intervention changes. This turns interaction reasoning into a queryable counterfactual object and supports branch-conditioned safety attribution.

\subsection{Risk-aware, social, and closed-loop planning}
Risk-aware and contingency planners optimize for safety under multi-modal futures or policy beliefs \parencite{mustafa2024racp}, and integrated prediction--planning systems learn costs or joint plans in interactive settings \parencite{huang2023differentiable,chen2023interactive}. Social and aggressiveness-aware planners encourage human-compatible behavior \parencite{crosato2024social,hu2023formulating}. Closed-loop resources such as WOMD, Waymax, WOSAC, nuPlan, and CARLA make it possible to evaluate interaction behavior beyond open-loop displacement \parencite{ettinger2021large,gulino2023waymax,montali2023waymo,caesar2021nuplan,dosovitskiy2017carla}. \framework{} is complementary to these lines: it does not replace collision risk, contingency planning, or closed-loop evaluation, but adds a mechanism-level criterion for whether a low-risk plan is non-coercive.

The mechanism-token architecture borrows standard machinery from conditional neural processes, attentive neural processes, Set Transformers, slot attention, and neural operators \parencite{garnelo2018conditional,kim2019attentive,lee2019set,locatello2020object,kovachki2023neural}. The novelty is not the attention module itself, but the driving-specific supervised object: same-root intervention groups, branch-conditioned burden and safety, forced-dependence labels, and planning through a calibrated mechanism-feasible set.

\begin{table}[t]
\centering
\caption{Conceptual distinction from representative related work. ``Surface'' denotes a same-root query function from ego interventions to response branch, burden, branch-conditioned safety, and motion.}
\resizebox{\textwidth}{!}{%
\begin{tabular}{lcccc}
\toprule
Method family & Ego-conditioned & Shared response surface & Burden / pressure & Forced-dependence planning \\
\midrule
PiP / DTPP / PPAD / HYPE & Yes & No or implicit & No & No \\
M2I / GameFormer & Partly & No or implicit & No & No \\
Behavior-latent / social-style models & No or partly & No & Partly & No \\
Risk-aware / contingency planners & Partly & No & Usually no & No \\
\framework{} & Yes & Yes & Yes & Yes \\
\bottomrule
\end{tabular}}
\label{tab:related_distinction}
\end{table}


\section{Method}
\label{sec:method}

\paragraph{Overview.}
We address a failure mode that is not captured by conventional collision-free planning: an ego trajectory can appear safe only because another road user is forced to absorb the conflict through a high-burden ceding response. We call such trajectories \emph{false-safe}. Our central idea is to replace ego-centric feasibility with \emph{non-coercive feasibility}: an ego plan should be considered feasible only if every critical road user still retains a sufficiently rich set of low-burden safe responses. This differs from courtesy-style objective shaping, where the effect on other agents is treated as an additional soft cost. In our formulation, safety-by-coercion is a feasibility defect.

Given a driving scene, our method first constructs a burden-oriented interaction graph that predicts not only who interacts with whom, but also who would have to absorb a potential conflict. It then proposes ego trajectory candidates, generates counterfactual natural alternatives for critical agents, evaluates ego-conditioned safe response sets, and rejects candidates with high-confidence coercion witnesses. The remaining candidates are ranked by standard ego utility.

\subsection{Burden-Oriented Interaction Graph}
\label{sec:method_graph}

Let the input scene be
\begin{equation}
    \mathcal{S} = (\mathcal{M}, \mathcal{H}, x_e),
\end{equation}
where $\mathcal{M}$ is the HD map, $\mathcal{H}=\{h_i\}_{i=1}^{N}$ denotes the historical states of surrounding agents, and $x_e$ is the current ego state. We encode $\mathcal{S}$ as a heterogeneous graph
\begin{equation}
    \mathcal{G}_B = (\mathcal{V}, \mathcal{E}_B),
\end{equation}
where nodes include the ego, surrounding agents, lane segments, conflict regions, and traffic-control elements. Unlike a standard interaction graph, each directed ego-agent edge is burden-oriented: it is designed to represent which agent would need to absorb a conflict, through which response channel, and with how much slack.

We introduce a conflict-query conditioned heterogeneous graph transformer. For a candidate ego-agent relation $(e,i)$, a compact conflict query $q_{ei}$ is computed from relative kinematics, map topology, and local right-of-way cues:
\begin{equation}
    q_{ei} = \phi_q(x_e, h_i, \mathcal{M}).
\end{equation}
The attention from node $u$ to node $v$ is modulated by this query:
\begin{equation}
    \alpha^{r}_{uv}
    =
    \operatorname{softmax}_{v}
    \left(
    \frac{
    (W^{r}_{Q} z_u)^\top
    (W^{r}_{K} z_v + W^{r}_{C} q_{uv})
    }{\sqrt{d}}
    \right),
\end{equation}
where $r$ is the heterogeneous relation type and $z_u,z_v$ are node embeddings. This attention has a physical interpretation: it asks which road user is likely to be forced to change behavior in a potential conflict.

The graph decoder predicts two edge states for each ego-agent pair:
\begin{equation}
    e^{\mathrm{nat}}_{ei}, \qquad e^{\mathrm{cond}}_{ei}.
\end{equation}
The natural edge $e^{\mathrm{nat}}_{ei}$ summarizes the agent's non-coerced passage relation, while the ego-conditioned edge $e^{\mathrm{cond}}_{ei}$ summarizes the response required after a specific ego candidate is considered. Their discrepancy is the first structural signal of false-safety:
\begin{equation}
    \Delta e_{ei}
    =
    \psi(e^{\mathrm{cond}}_{ei})-\psi(e^{\mathrm{nat}}_{ei}),
\end{equation}
where $\psi(\cdot)$ projects an edge state to conflict timing, priority, slack, response type, and burden-relevant attributes.

The decoder also outputs a set of candidate coercion witnesses
\begin{equation}
    \widehat{\mathcal{W}} =
    \{ \widehat{W}_{ei} \}_{i=1}^{N},
\end{equation}
where each witness proposes a burdened agent, a conflict interval, a ceding mechanism token, and a burden magnitude. These witnesses are later validated by counterfactual rollout rather than trusted as pure neural predictions.

We train the graph and witness decoder with a compact witness loss:
\begin{equation}
    \mathcal{L}_{\mathrm{wit}}
    =
    \lambda_w \mathcal{L}_{\mathrm{exist}}
    +
    \lambda_z \mathcal{L}_{\mathrm{token}}
    +
    \lambda_b \mathcal{L}_{\mathrm{burden}}
    +
    \lambda_c \mathcal{L}_{\mathrm{conflict}} .
\end{equation}
Here $\mathcal{L}_{\mathrm{exist}}$ supervises whether an ego-agent pair contains a coercion witness, $\mathcal{L}_{\mathrm{token}}$ classifies the burden mechanism, $\mathcal{L}_{\mathrm{burden}}$ regresses burden severity, and $\mathcal{L}_{\mathrm{conflict}}$ localizes the conflict interval. The construction of pseudo-labels is detailed in Appendix~\ref{app:witness_supervision}.

\subsection{Ego Candidate Generation and Non-Coercive Feasibility}
\label{sec:method_candidate}

The ego candidate generator takes the scene and the burden-oriented graph as input:
\begin{equation}
    \{(\tau^k_e, a^k_e)\}_{k=1}^{K}
    =
    \Pi_{\theta}(\mathcal{S}, \mathcal{G}_B),
\end{equation}
where $\tau^k_e$ is a continuous ego trajectory and $a^k_e$ is a proposal descriptor. In the evaluated implementation, the bank is a finite set of deterministic kinematic primitives: constant-acceleration, comfort-bounded stop/creep, conflict-timing, terminal-state, and delayed lateral-offset trajectories. Synthetic primitives are screened for acceleration, jerk, duplicate endpoints, and sampled lane-corridor compliance. The macro label describes how a primitive was generated; it is not itself evidence of route feasibility, right-of-way, or successful negotiation.

For each ego candidate $\tau^k_e$, let $\mathcal{I}^k$ be the set of critical agents selected from $\mathcal{G}_B$. We define the protected set
\begin{equation}
    \mathcal{P}^k
    =
    \left\{i\in\mathcal{I}^k:
    \rho^k_{ei}\in\{\mathrm{AgentPriority},\mathrm{EqualOrNegotiated}\}
    \right\},
    \label{eq:protected_set}
\end{equation}
using rule-derived priority as an anchor and a learned priority score only as a calibrated refinement. Hard non-coercive feasibility is defined on $\mathcal{P}^k$; the all-critical set $\mathcal{I}^k$ is retained for global burden-transfer diagnostics. For each $i\in\mathcal{I}^k$, let $\mathcal{A}^{\mathrm{nat}}_i$ denote the non-coerced natural alternative set of agent $i$, and let $\mathcal{R}^{\mathrm{safe}}_i(\tau^k_e)$ denote the set of ego-conditioned safe responses. The low-burden safe response set is
\begin{equation}
    \mathcal{R}^{\mathrm{low}}_i(\tau^k_e)
    =
    \left\{
    \tau_i \in \mathcal{R}^{\mathrm{safe}}_i(\tau^k_e)
    \mid
    B_i(\tau_i;\mathcal{S},\tau^k_e)
    \le
    \beta_i(\mathcal{S})
    \right\},
\end{equation}
where $B_i$ is a hybrid burden functional and $\beta_i$ is a scene-conditioned burden threshold.

We represent the natural alternative set by stable roots $r_{im}$ with calibrated masses $p_{im}$, $\sum_m p_{im}=1$. For root $m$, let $c_{ikm}\in\{0,1\}$ indicate whether the natural trajectory conflicts with ego candidate $k$, and let $q_{ikm}\in[0,1]$ denote the probability that a same-root, dynamically valid, low-burden safe response exists. The retained-root probability is
\begin{equation}
    s_{ikm}=(1-c_{ikm})\,r_{ikm}+c_{ikm}\,q_{ikm},
\end{equation}
where $r_{ikm}$ denotes low-burden retention without a conflict. The option-preservation ratio is the transported natural mass
\begin{equation}
    O_i(\tau^k_e)=\sum_m \widetilde p_{im}s_{ikm},
    \qquad
    \widetilde p_{im}=(1-\varepsilon_p)p_{im}+\varepsilon_p/M,
    \label{eq:option_preservation}
\end{equation}
where the small mass floor $\varepsilon_p$ prevents a learned predictor from hiding a safety-relevant root by assigning it zero probability. Before floor smoothing, roots with calibrated mass below $p_{\min}$ are removed and the surviving mass is renormalized; label generation, training adaptation, and inference use this same measure. A candidate is priority-aware non-coercively feasible only if every $i\in\mathcal{P}^k$ has at least one same-root low-burden safe response and $O_i(\tau^k_e)\ge\alpha$. The corresponding all-critical condition is reported as a global stress diagnostic rather than used as the primary hard veto. This finite, root-indexed definition replaces an ill-defined set intersection between continuous trajectory samples and matches the implemented transport labels.

\paragraph{Root-conditioned recovery--burden representation.}
For every candidate--agent--root tuple, the transport head predicts $q_{ikm}$ and $b^\star_{ikm}$ from a shared root-conditioned latent but through separate output channels.  This separation is necessary: $q_{ikm}=0$ means that no safe response lies below $\beta_i$, whereas a finite same-root safe response with $b^\star_{ikm}>\beta_i$ may still exist.  We add a smooth consistency term between $q_{ikm}$ and $\sigma((\beta_i-b^\star_{ikm})/T_b)$ on confidently labelled conflicting roots, while retaining direct supervision for both quantities.  A generic compact response bank is decoded only as an auxiliary reconstruction and visualization path; it does not define root recovery or the primary burden-tail certificate.

\subsection{Counterfactual Natural Alternatives}
\label{sec:method_counterfactual}

A standard predictor may learn that surrounding agents often yield to the ego vehicle, which can contaminate the natural alternative set. To avoid treating coerced yielding as natural behavior, we construct $\mathcal{A}^{\mathrm{nat}}_i$ through three complementary branches:
\begin{equation}
    \mathcal{A}^{\mathrm{nat}}_i
    =
    \mathcal{A}^{\mathrm{obs}}_i
    \cup
    \mathcal{A}^{\mathrm{neu}}_i
    \cup
    \mathcal{A}^{\mathrm{prio}}_i .
    \label{eq:natural_union}
\end{equation}

The observational branch $\mathcal{A}^{\mathrm{obs}}_i$ uses a multimodal trajectory decoder conditioned on the full scene to keep the alternatives close to the empirical driving distribution. The ego-neutral branch $\mathcal{A}^{\mathrm{neu}}_i$ predicts the agent's response under a \emph{model-based intervention proxy} in which candidate-specific ego pressure is removed:
\begin{equation}
    \mathcal{A}^{\mathrm{neu}}_i
    \sim
    p_\theta\!\left(\tau_i \mid \operatorname{Neutralize}(\mathcal{S},\tau_e^k),\tau^{\mathrm{neu}}_e\right).
\end{equation}
We deliberately avoid interpreting this distribution as an identified causal effect from logged data alone; its validity is evaluated with independent reactive and human-audited protocols. The priority-preserving branch $\mathcal{A}^{\mathrm{prio}}_i$ generates rule-constrained alternatives that preserve the agent's local priority and comfort-bounded motion.

Each source is represented by stable typed roots and a causal dynamics residual. For root $m$, the decoder predicts bounded local longitudinal/lateral acceleration, jerk, and yaw-rate corrections $\delta u_{im}(t)$ around an analytic source prototype, and obtains position and velocity by integration. Source-adaptive control capacity gives observational roots enough freedom to represent real curvature while neutral and priority-preserving roots remain more conservative.

A fixed endpoint ball is not an adequate definition of root identity: it treats a low-probability exploratory root and a probability-carrying root equally, ignores large intermediate excursions that return near the endpoint, and can suppress a valid observational maneuver. Because OPR is defined on retained probability mass, COWP instead uses a probability-mass-aware, multi-horizon semantic envelope. Let
\begin{equation}
    \rho_{im}
    =
    \max_t
    \frac{\|\widetilde{\delta p}_{im}(t)\|_2}{r_{s(m)}(t)}
\end{equation}
be the maximum pre-projection path ratio for root $m$, and let $\widetilde p_{im}=(1-\varepsilon_p)p_{im}+\varepsilon_p/M$ be the same floor-smoothed root mass used by OPR. The implemented identity loss is
\begin{equation}
    \mathcal{L}_{\mathrm{id}}
    =
    \frac{
      \sum_{i,m}\operatorname{sg}(\widetilde p_{im})\,
      \mathbf{1}[i\in\mathcal I]
      \left[
        \frac{\rho_{im}-\eta}{1-\eta}
      \right]_+^2
    }{
      \sum_{i,m}\operatorname{sg}(\widetilde p_{im})\,
      \mathbf{1}[i\in\mathcal I]
    },
    \label{eq:natural_mass_envelope}
\end{equation}
where $\eta<1$ defines an interior margin and $\operatorname{sg}$ denotes stop-gradient. Detaching the mass prevents the decoder from evading geometric regularization by suppressing a violating root's logit, while the probability floor prevents it from making that root completely invisible. Separately, every mode is constrained by a wider emergency envelope $\bar r_{s(m)}(t)>r_{s(m)}(t)$. If a raw integrated residual exceeds this envelope, the complete control sequence is scaled and re-integrated. Thus all modes remain physically bounded, while semantic root-identity pressure is concentrated on the probability mass that enters OPR.

After generation, all alternatives are filtered by map compliance, dynamic feasibility, comfort bounds, and plausibility. The filtered set is used as the reference against which ego-conditioned coercion is evaluated. The natural alternative generator is trained with
\begin{equation}
\begin{aligned}
    \mathcal{L}_{\mathrm{nat}}
    ={}&
    \mathcal{L}_{\mathrm{set}}
    +\lambda_{\mathrm{prior}}\mathcal{L}_{\mathrm{prior}}
    +\lambda_{\mathrm{ctrl}}\mathcal{L}_{\mathrm{ctrl}}
    +\lambda_{\mathrm{mode}}\mathcal{L}_{\mathrm{usage}}
    +\lambda_{\mathrm{id}}\mathcal{L}_{\mathrm{id}}.
\end{aligned}
\end{equation}
Here $\mathcal{L}_{\mathrm{set}}$ contains source-restricted weighted set prediction and short-horizon terms; $\mathcal{L}_{\mathrm{prior}}$ softly protects neutral and priority-preserving typed roots; $\mathcal{L}_{\mathrm{ctrl}}$ regularizes control variation; $\mathcal{L}_{\mathrm{usage}}$ discourages root collapse; and $\mathcal{L}_{\mathrm{id}}$ is Eq.~\eqref{eq:natural_mass_envelope}. We do not add a separate observational ``gain'' loss to the main objective: observational accuracy is already optimized by the source-restricted set loss, and an additional basis-relative shortfall term is accepted only if an isolated ablation demonstrates incremental benefit. Component claims are made only under identical initialization, data order, encoder state, precision policy, and checkpoint-selection protocol. Detailed construction is provided in Appendix~\ref{app:natural_alternatives}.

\subsection{Coercion-Witness Certification and Plan Selection}
\label{sec:method_certification}

For each ego candidate $\tau^k_e$ and critical agent $i$, we roll out the candidate against root-indexed natural alternatives and ego-conditioned responses. Let $b^{\star}_{ikm}$ be the minimum primitive burden among safe responses assigned to natural root $m$, with $b^{\star}_{ikm}=+\infty$ when no safe response exists. We define a tail burden excess over conflicted roots,
\begin{equation}
    C_i(\tau^k_e)
    =\operatorname{CVaR}_{\rho}
    \left(
      \left\{[b^{\star}_{ikm}-\beta_i]_+\right\}_{m:c_{ikm}=1};
      \widetilde p_{im}
    \right),
\end{equation}
which avoids subtracting minima from unrelated roots. Let $M^{\mathrm{conf}}_{ik}=\sum_m\widetilde p_{im}c_{ikm}$ be natural conflict mass. A coercion witness is confirmed only when a non-negligible natural mass is blocked and the best same-root escape is high burden or unavailable:
\begin{equation}
    W^k_{ei}=\mathbf{1}
    \left[
      M^{\mathrm{conf}}_{ik}>\delta_c
      \ \wedge\
      \left(C_i(\tau^k_e)>\gamma\ \vee\ O_i(\tau^k_e)<\alpha\right)
    \right].
    \label{eq:witness}
\end{equation}
The earlier existential condition ``there exists a high-burden safe response'' is insufficient because nearly every scene admits an emergency maneuver; the certificate must depend on the minimum or tail of the best safe responses.

To compress root-level evidence without replacing the mechanism by an unconstrained candidate classifier, we use a monotone Budgeted Counterfactual Option Transport (BCOT) head. For each protected pair, the deficit combines unrecovered conflicted mass, conflict-conditioned burden excess, and normalized OPR shortfall. Candidate risk aggregates a priority-weighted mean, a smooth upper tail, a severe-pair term, and the maximum pair deficit. All aggregation weights are constrained positive and normalized. A second all-critical BCOT output is trained and reported as a global diagnostic, but only the protected-priority output defines the primary hard certificate. This two-level construction matches the paper's priority-aware claim while preserving visibility into burden transferred to non-protected agents.

Because natural roots form an unordered set, per-root transport supervision requires an alignment between generated roots and cached natural alternatives. We align only for supervision, without gradient, using a shared multi-horizon geometric cost over 1/3/5/8~s plus a penalty for crossing the structural source identity (observational, neutral, or priority-preserving). This prevents a geometrically similar yielding root from being relabeled as the same observational root after an intervention. The alignment is diagnostic and does not change the retained-mass definition itself.

During inference, safety and non-coercive feasibility are applied in a hard-first manner:
\begin{equation}
    \mathcal{K}_{\mathrm{safe}}
    =
    \{k \mid S_{\mathrm{coll}}(\tau^k_e)\le \epsilon \},
\end{equation}
\begin{equation}
    \mathcal{K}_{\mathrm{ncf}}
    =
    \{k\in\mathcal{K}_{\mathrm{safe}}
    \mid
    C_i(\tau^k_e)\le \gamma,\ 
    O_i(\tau^k_e)\ge \alpha,\ 
    \forall i\in\mathcal{P}^k
    \}.
\end{equation}
High-confidence coercion witnesses cause rejection; uncertain witnesses are retained with an additional penalty. The final selected trajectory is
\begin{equation}
    \tau^\star_e
    =
    \arg\min_{k\in\mathcal{K}_{\mathrm{ncf}}}
    J_{\mathrm{ego}}(\tau^k_e),
\end{equation}
where $J_{\mathrm{ego}}$ includes progress, comfort, route following, and lane regularity. If $\mathcal{K}_{\mathrm{ncf}}$ is empty, the planner executes a conservative fallback policy.

Training uses soft surrogates for differentiability:
\begin{equation}
    \mathcal{L}
    =
    \mathcal{L}_{\mathrm{wit}}
    +
    \mathcal{L}_{\mathrm{nat}}
    +
    \mathcal{L}_{\mathrm{resp}}
    +
    \mathcal{L}_{\mathrm{plan}},
\end{equation}
where $\mathcal{L}_{\mathrm{resp}}$ supervises ego-conditioned safe response classification and low-burden response estimation, and $\mathcal{L}_{\mathrm{plan}}$ combines imitation, non-coercive ranking, and closed-loop outcome losses. The complete burden function, safe response construction, fallback policy, and training targets are given in Appendix~\ref{app:implementation_details}.



\section{Experiments}
\label{sec:experiments}

\subsection{Experimental Setup}
\label{sec:exp_setup}

We evaluate COWP on interaction-heavy driving scenarios where conventional collision-free planning may hide burden transfer to other road users. Our primary evaluation uses the Waymo Open Motion Dataset and Waymax simulator, which provide real-world multi-agent trajectories, HD maps, and closed-loop multi-agent simulation support \cite{ettinger2021large, gulino2023waymax}. We focus on scenarios containing lane changes, merges, unprotected turns, intersection crossings, and dense car-following, where ego decisions can plausibly force another vehicle to brake, yield, or surrender a gap.

For each scenario, we initialize the ego and surrounding agents from logged states and evaluate the ego planner in closed loop. The primary Waymax protocol uses honest logged replay for non-ego agents; it supports conventional SDC collision, off-road, wrong-way, and route-progress measurements, but it does not provide ground-truth counterfactual burden. We therefore report learned burden quantities as mechanism diagnostics in this protocol. The final causal burden claim additionally requires a separately trained reactive-agent protocol and a held-out human-audited false-safe stress set; results from these protocols are never pooled with logged replay.

\paragraph{Evidence protocol.}
An engineering smoke run validates only data flow and numerical integrity. Decoder, loss, and residual-capacity claims require fresh one-factor ablations with identical initialization, data order, encoder state, precision policy, evaluator, sampled scene IDs, and checkpoint epoch. A single-seed natural-component gate is used only to decide whether downstream evidence collection is justified. Publication-level decoder, planner, and selector claims require at least three independent seeds, paired scenario-level confidence intervals, no bypassed quality gate, and a natural-basis checkpoint that has already passed its protocol-aligned gate. Downstream results are treated as mechanism probes rather than evidence when the natural basis is unverified.

\paragraph{Baselines.}
We compare COWP against two groups of baselines. The mechanism-controlled group shares the same candidate bank and model checkpoint: conventional physical-safety selection, planner-score-only selection, soft-burden ranking, pairwise-max witness rejection, and an epsilon-Pareto selector. This group isolates the contribution of root transport and hard non-coercive feasibility. The external planner group includes reproducible implementations of IDM/Lattice, PDM-Closed, PlanT, and GameFormer where their input and route-conditioning assumptions can be matched. All external comparisons use the same scenario IDs, horizon, action projection, and Waymax metric adapter.

\subsection{Evaluation Metrics}
\label{sec:metrics}

Standard planning metrics such as collision rate and progress are necessary but insufficient for evaluating false-safe behavior. We therefore report both conventional metrics and non-coercive interaction metrics.

\paragraph{Collision Rate (CR), Off-road Rate (OR), and Closed-loop Unsafe Rate (CUR).}
CR counts ego--agent overlap and OR counts leaving the drivable area; CUR is their union:
\begin{equation}
    \mathrm{CR}=\frac{1}{N}\sum_{n=1}^{N}\mathbf{1}[\mathrm{Collision}_n],
    \qquad
    \mathrm{OR}=\frac{1}{N}\sum_{n=1}^{N}\mathbf{1}[\mathrm{OffRoad}_n],
\end{equation}
\begin{equation}
    \mathrm{CUR}=\frac{1}{N}\sum_n\mathbf{1}[\mathrm{Collision}_n\vee\mathrm{OffRoad}_n].
\end{equation}

\paragraph{Ego Progress (EP).}
The route-normalized ego progress:
\begin{equation}
    \mathrm{EP}
    =
    \frac{1}{N}
    \sum_{n=1}^{N}
    \frac{
    s^{n}_{e}(T)-s^{n}_{e}(0)
    }{
    s^{n}_{\mathrm{route}}
    }.
\end{equation}
This metric verifies that reducing coercion does not simply produce an overly conservative planner.

\paragraph{Priority Burden-Transfer Rate (PBTR).}
PBTR is the primary task-specific metric. Among selected plans that are conventionally collision-free and contain at least one protected relation, it measures how often the plan still creates a coercion witness for a protected agent:
\begin{equation}
    \mathrm{PBTR}
    =
    \frac{
      \sum_n \mathbf 1[
      \mathrm{CollisionFree}_n
      \wedge |\mathcal P^n|>0
      \wedge \exists i\in\mathcal P^n:W^n_{ei}=1]
    }{
      \sum_n \mathbf 1[
      \mathrm{CollisionFree}_n
      \wedge |\mathcal P^n|>0]
    }.
\end{equation}
PBTR most directly tests the paper's priority-aware claim: ego safety should not be obtained by transferring conflict-resolution burden to an agent whose priority must be protected or negotiated.

\paragraph{Global False-Safe Rate (FSR).}
For transparency, we also report the stricter all-critical diagnostic
\begin{equation}
    \mathrm{FSR}
    =
    \frac{
    \sum_{n}
    \mathbf{1}
    [
    \mathrm{CollisionFree}_n
    \wedge
    \exists i\in\mathcal I^n: W^n_{ei}=1
    ]
    }{
    \sum_{n}
    \mathbf{1}[\mathrm{CollisionFree}_n]
    }.
\end{equation}
FSR is intentionally broader than the hard certificate and reveals whether a priority-aware planner merely moves burden to non-protected agents.

\paragraph{Protected Option Preservation Ratio (OPR).}
OPR is the retained probability mass of same-root low-burden responses, averaged over protected relations:
\begin{equation}
    \mathrm{OPR}_{\mathcal P}
    =
    \frac{1}{N}
    \sum_{n=1}^{N}
    \frac{1}{|\mathcal{P}^n|}
    \sum_{i\in\mathcal{P}^n}
    \sum_m \widetilde p^n_{im}s^n_{im}.
\end{equation}
We report the all-critical OPR only as an auxiliary diagnostic. OPR is not interpreted in isolation: a method can preserve options by stopping excessively, so it is paired with progress and coverage.

\paragraph{Protected Tail Burden-Transfer Excess (BTE-CVaR$_{25}$).}
PBTR measures frequency but not severity.  For a selected plan in scene $n$, we first compute the root-conditioned, probability-weighted burden-excess tail for every protected relation and retain the worst protected relation,
\begin{equation}
    e_n
    =
    \max_{i\in\mathcal P^n}
    \operatorname{CVaR}_{\rho}
    \left(
      \{[b^{\star,n}_{im}-\beta_i^n]_+\}_{m:c^n_{im}=1};
      \widetilde p^n_{im}
    \right).
\end{equation}
We then report the empirical mean of the worst quartile of $\{e_n\}$,
\begin{equation}
    \mathrm{BTE\mbox{-}CVaR}_{25}
    =
    \operatorname{CVaR}_{25\%}\left(\{e_n\}_{n:\,|\mathcal P^n|>0}\right).
\end{equation}
This two-level construction matches the implementation: root uncertainty is summarized within each protected relation, the most burdened protected agent represents the selected plan, and the outer tail prevents a small number of severe transfers from being hidden by an average.  CBS and HBCR are retained only as appendix diagnostics for continuity with earlier versions.

\paragraph{NCF Scene Retention and Non-Coercive Progress Regret.}
Proposal coverage and certificate conservatism are separated by
\begin{equation}
    \mathrm{NCF\mbox{-}Ret}
    =
    \frac{
      \sum_n \mathbf 1[\mathcal K^n_{\mathrm{prop,ncf}}\neq\emptyset
      \wedge \mathcal K^n_{\mathrm{cert,ncf}}\neq\emptyset]
    }{
      \sum_n \mathbf 1[\mathcal K^n_{\mathrm{prop,ncf}}\neq\emptyset]
    },
\end{equation}
and, on scenes with at least one non-coercive proposal,
\begin{equation}
    \mathrm{NPR}
    =
    \frac{1}{N_{\mathrm{ncf}}}
    \sum_n
    \frac{[s^{n}_{\mathrm{best\mbox{-}ncf}}-s^{n}_{\mathrm{selected}}]_+}
         {\max(s^{n}_{\mathrm{best\mbox{-}ncf}},\epsilon)}.
\end{equation}
NCF-Ret diagnoses certificate false rejection, while NPR detects whether the selector keeps a feasible candidate but ranks an unnecessarily conservative one. We additionally report a PBTR--coverage curve rather than relying on a single calibrated budget.

\paragraph{Witness Localization Accuracy (WLA).}
For scenarios with pseudo-label or human-verified witnesses, we evaluate whether the method identifies the correct burdened agent and conflict interval:
\begin{equation}
    \mathrm{WLA}
    =
    \frac{1}{|\mathcal{D}_{\mathrm{wit}}|}
    \sum_{(n,i)\in\mathcal{D}_{\mathrm{wit}}}
    \mathbf{1}
    [
    \widehat{i}=i
    \wedge
    \mathrm{IoU}(\widehat{T}_{ei},T_{ei})>\delta_T
    ].
\end{equation}

\paragraph{Mechanism Token Accuracy (MTA).}
For positive witness cases, we measure whether the predicted mechanism token matches the pseudo-label:
\begin{equation}
    \mathrm{MTA}
    =
    \frac{1}{|\mathcal{D}_{\mathrm{wit}}|}
    \sum_{(n,i)\in\mathcal{D}_{\mathrm{wit}}}
    \mathbf{1}
    [
    \widehat{z}_{ei}=z^{\mathrm{gt}}_{ei}
    ].
\end{equation}
This evaluates whether the method distinguishes hard braking, abrupt yielding, priority abandonment, gap surrender, safety-reserve consumption, and option reduction.

\subsection{Main Results}
\label{sec:main_results}

Table~\ref{tab:main_results} reports closed-loop results on interaction-heavy scenarios. The main comparison is not only whether a planner avoids collision, but whether it avoids doing so by transferring burden to surrounding agents.

\begin{table}[t]
\centering
\caption{Closed-loop planning performance on interaction-heavy scenarios. Standard safety is reported with CR/OR; PBTR is the primary mechanism metric. OPR and NCF-Ret must be interpreted jointly with EP.}
\label{tab:main_results}
\resizebox{\linewidth}{!}{
\begin{tabular}{lcccccc}
\toprule
Method & CR $\downarrow$ & OR $\downarrow$ & EP $\uparrow$ & PBTR $\downarrow$ & OPR$_{\mathcal P}$ $\uparrow$ & NCF-Ret $\uparrow$ \\
\midrule
IDM/Lattice & -- & -- & -- & -- & -- & -- \\
PDM-Closed \cite{dauner2023parting} & -- & -- & -- & -- & -- & -- \\
PlanT \cite{renz2022plant} & -- & -- & -- & -- & -- & -- \\
GameFormer \cite{gameformer} & -- & -- & -- & -- & -- & -- \\
\midrule
COWP w/o Counterfactual Alternatives & -- & -- & -- & -- & -- & -- \\
COWP w/o Option Preservation & -- & -- & -- & -- & -- & -- \\
COWP w/o Witness Rejection & -- & -- & -- & -- & -- & -- \\
\textbf{COWP} & -- & -- & -- & -- & -- & -- \\
\bottomrule
\end{tabular}
}
\end{table}

Conventional planners may remain competitive on collision rate and ego progress while exhibiting higher priority burden transfer in dense interactions. COWP is evaluated primarily by whether it reduces PBTR and BTE-CVaR$_{25}$ at comparable CR, OR, and EP, while retaining non-coercive proposal coverage.

\subsection{False-Safe Stress Test}
\label{sec:false_safe_stress}

To directly test the proposed failure mode, we construct a false-safe stress set from logged interaction scenarios. For each original scenario, we generate ego intervention candidates that vary merge timing, crossing priority, target gap selection, and longitudinal aggressiveness. These interventions include both non-coercive candidates and candidates that remain collision-free only if another agent yields with high burden.

\begin{table}[t]
\centering
\caption{False-safe stress test. Accept Rate is computed separately for non-coercive and coercive candidates. A desirable planner accepts non-coercive candidates and rejects coercive collision-free candidates.}
\label{tab:stress}
\resizebox{\linewidth}{!}{
\begin{tabular}{lcccc}
\toprule
Method & Accept Non-Coercive $\uparrow$ & Accept False-Safe $\downarrow$ & Witness Recall $\uparrow$ & Witness Precision $\uparrow$ \\
\midrule
IDM/Lattice & -- & -- & -- & -- \\
PDM-Closed \cite{dauner2023parting} & -- & -- & -- & -- \\
PlanT \cite{renz2022plant} & -- & -- & -- & -- \\
GameFormer \cite{gameformer} & -- & -- & -- & -- \\
\textbf{COWP} & -- & -- & -- & -- \\
\bottomrule
\end{tabular}
}
\end{table}

This test isolates the key claim of the paper: a method should not merely predict that another vehicle will yield; it should determine whether relying on that yield is non-coercively feasible.

\subsection{Ablation Study}
\label{sec:ablation}

Table~\ref{tab:ablation} evaluates the contribution of each component. Removing counterfactual alternatives tests whether standard prediction is contaminated by observed yielding behavior. Removing option preservation tests whether existence of one low-burden response is sufficient. Removing witness rejection tests whether burden should be a soft penalty or a feasibility-level certificate.

\begin{table}[t]
\centering
\caption{Ablation study of COWP components.}
\label{tab:ablation}
\resizebox{\linewidth}{!}{
\begin{tabular}{lcccc}
\toprule
Variant & PBTR $\downarrow$ & BTE-CVaR$_{25}$ $\downarrow$ & OPR$_{\mathcal P}$ $\uparrow$ & NCF-Ret $\uparrow$ \\
\midrule
Full COWP & -- & -- & -- & -- \\
w/o Ego-Neutral Branch & -- & -- & -- & -- \\
w/o Priority-Preserving Branch & -- & -- & -- & -- \\
w/o Option Preservation & -- & -- & -- & -- \\
w/o Witness Rejection & -- & -- & -- & -- \\
Soft Burden Cost Only & -- & -- & -- & -- \\
\bottomrule
\end{tabular}
}
\end{table}

The key test is whether a hard protected-priority certificate reduces PBTR beyond a soft burden cost without collapsing NCF-Ret or ego progress. This trend is reported only after the held-out mechanism and online Waymax gates pass.

\subsection{Witness Quality and Interpretability}
\label{sec:witness_quality}

We evaluate whether COWP identifies the correct burdened agent, conflict interval, and ceding mechanism. Table~\ref{tab:witness} reports witness localization and token classification.

\begin{table}[t]
\centering
\caption{Coercion witness quality.}
\label{tab:witness}
\resizebox{\linewidth}{!}{
\begin{tabular}{lcccccc}
\toprule
Method & WLA $\uparrow$ & MTA $\uparrow$ & HB-F1 $\uparrow$ & AY-F1 $\uparrow$ & PA-F1 $\uparrow$ & GS-F1 $\uparrow$ \\
\midrule
COWP w/o Dual Edge & -- & -- & -- & -- & -- & -- \\
COWP w/o Conflict Query & -- & -- & -- & -- & -- & -- \\
\textbf{COWP} & -- & -- & -- & -- & -- & -- \\
\bottomrule
\end{tabular}
}
\end{table}

Qualitative results will visualize the accepted ego trajectory, rejected false-safe candidates, the burdened agent, the natural alternative, the required ceding response, and the mechanism token. These visualizations are used to test whether COWP rejects plans because they transfer conflict-resolution responsibility to another road user rather than merely because they are slow or conservative.

\section{Conclusion}
\label{sec:conclusion}

We studied a false-safe failure mode in interactive autonomous driving: an ego trajectory can be collision-free only because another road user is forced to absorb the conflict through a high-burden ceding response. This failure mode is difficult to capture with standard collision, comfort, or route-following metrics, and it is not fully addressed by treating social behavior as a soft planning cost.

We proposed COWP, a coercion-witness planning framework based on non-coercive feasibility. COWP evaluates whether each critical road user retains a sufficiently rich set of low-burden safe responses under a candidate ego plan. It constructs burden-oriented interaction graphs, generates counterfactual natural alternatives, evaluates ego-conditioned safe response sets, and rejects collision-free candidates when a coercion witness shows that safety depends on another agent's costly yielding behavior.

Our formulation reframes interactive safety from ``Can the ego avoid collision?'' to ``Can the ego remain safe without forcing others to absorb the conflict?'' This perspective exposes a class of planning failures that conventional benchmarks can overlook. Future work will extend non-coercive feasibility to pedestrian and cyclist interactions, richer social conventions, and uncertainty-calibrated certificates in large-scale closed-loop deployment.

\appendix

\section{Implementation Details}
\label{app:implementation_details}

\subsection{Graph Construction and Edge Semantics}
\label{app:graph_details}

We construct a heterogeneous graph with five node types:
\begin{equation}
    \mathcal{V}
    =
    \mathcal{V}_{\mathrm{ego}}
    \cup
    \mathcal{V}_{\mathrm{agent}}
    \cup
    \mathcal{V}_{\mathrm{lane}}
    \cup
    \mathcal{V}_{\mathrm{conflict}}
    \cup
    \mathcal{V}_{\mathrm{ctrl}} .
\end{equation}
Agent nodes store position, velocity, acceleration, heading, yaw rate, size, semantic type, and history embedding. Lane nodes store centerline geometry, speed limit, turn direction, and lane connectivity. Conflict nodes are created from lane-lane intersections, merge areas, crossing points, and predicted swept-volume overlaps. Control nodes represent traffic lights, stop signs, yield signs, and priority cues.

For an ego-agent pair $(e,i)$, the burden-oriented edge state is represented as
\begin{equation}
    e_{ei}
    =
    (
    c_{ei},
    \Delta t_{ei},
    \rho_{ei},
    r_{ei},
    s_{ei},
    \widehat{b}_{ei},
    z_{ei}
    ),
\end{equation}
where $c_{ei}$ is the conflict region, $\Delta t_{ei}$ is the temporal overlap around the conflict, $\rho_{ei}$ is the priority relation, $r_{ei}$ is the required response type, $s_{ei}$ is available slack, $\widehat{b}_{ei}$ is the predicted burden prior, and $z_{ei}$ is a mechanism token. In the main paper we avoid expanding this tuple to reduce notation load, but it is used explicitly in implementation.

\paragraph{Conflict query.}
The compact query $q_{ei}$ is computed as
\begin{equation}
    q_{ei}
    =
    \mathrm{MLP}
    \left[
    \Delta p_{ei},
    \Delta v_{ei},
    \Delta \psi_{ei},
    \Delta \mathrm{TTA}_{ei},
    d^{\mathrm{conf}}_{e},
    d^{\mathrm{conf}}_{i},
    \rho_{ei}
    \right],
\end{equation}
where $\Delta \mathrm{TTA}_{ei}$ is the difference between ego and agent time-to-arrival at the nearest conflict region, and $d^{\mathrm{conf}}_e,d^{\mathrm{conf}}_i$ are distances to that region.

\paragraph{Dual-edge decoding.}
The graph decoder predicts
\begin{equation}
    e^{\mathrm{nat}}_{ei}
    =
    D_{\mathrm{nat}}(z_e,z_i,q_{ei}),
    \qquad
    e^{\mathrm{cond}}_{ei}
    =
    D_{\mathrm{cond}}(z_e,z_i,q_{ei},\tau^k_e).
\end{equation}
The first edge is decoded without the candidate ego trajectory and therefore estimates the non-coerced passage relation. The second edge is candidate-conditioned and estimates what the other agent must do if $\tau^k_e$ is executed.

\paragraph{Mechanism tokens.}
We use six burden mechanism tokens:
\begin{equation}
    z_{ei}\in
    \{
    \mathrm{HB},
    \mathrm{AY},
    \mathrm{PA},
    \mathrm{GS},
    \mathrm{SR},
    \mathrm{OR}
    \}.
\end{equation}
They denote hard braking, abrupt yielding, priority abandonment, gap surrender, safety-reserve consumption, and option reduction. Tokens are used for intermediate supervision, selection constraints, and post-hoc explanations.

\subsection{Proposal Descriptors and Ego Candidate Generation}
\label{app:ego_candidates}

Each ego candidate is represented by a proposal descriptor and a continuous trajectory:
\begin{equation}
    (\tau^k_e,a^k_e), \qquad
    \tau^k_e=\{x^k_e(t)\}_{t=1}^{T}.
\end{equation}
The descriptor set is
\begin{equation}
\begin{aligned}
    \mathcal{A}_{\mathrm{prop}}
    =
    \{
    &\mathrm{KeepLane},
    \mathrm{Yield},
    \mathrm{Creep},
    \mathrm{StopBeforeConflict},\\
    &\mathrm{MergeAhead},
    \mathrm{MergeBehind},
    \mathrm{LaneChangeLeft},
    \mathrm{LaneChangeRight},\\
    &\mathrm{AccelerateCross},
    \mathrm{DecelerateCross}
    \}.
\end{aligned}
\end{equation}
These labels identify the primitive family and support stratified analysis; they do not certify that a lane change exists, that a gap is negotiable, or that a maneuver is map-feasible.

The evaluated proposal bank contains: logged ego replay; constant-acceleration longitudinal primitives; smooth stop and creep primitives; conflict-region stopping margins; conflict-arrival timing offsets; delayed lateral-offset primitives; and terminal speed/position variants.  We further introduce a \emph{Bidirectional Conflict-Time Envelope} (BCTE) proposal family.  For a conflict region at longitudinal distance $d$ and a nearby agent with estimated arrival time $t_j$, BCTE proposes ego arrival times $t_j-\Delta$ and $t_j+\Delta$ and solves the bounded constant-acceleration primitive
\begin{equation}
    a_e(d,t)=\frac{2(d-v_{e,0}t)}{t^2},
    \qquad
    a_{\min}\le a_e(d,t)\le a_{\max}.
\end{equation}
The two sides explicitly represent pass-before and pass-after hypotheses.  BCTE changes only proposal coverage: every generated trajectory remains subject to conventional safety and the root-conditioned non-coercive certificate, so adding proposals cannot turn an uncertified maneuver into a certified one.

For each generated trajectory we enforce finite state, acceleration, deceleration, jerk, and endpoint-diversity checks. We additionally screen sampled points against a locally cropped WOMD lane-centerline cloud:
\begin{equation}
    \frac{1}{|\mathcal T_s|}
    \sum_{t\in\mathcal T_s}
    \mathbf 1[d(x_e(t),\mathcal M_{\mathrm{lane}})\le d_{\max}]
    \ge \zeta,
    \qquad
    \max_{t\in\mathcal T_s}d(x_e(t),\mathcal M_{\mathrm{lane}})\le \bar d.
\end{equation}
This is a fast proposal-level map screen, not a claim of route-conditioned lattice-MPC optimality. Candidates are retained up to the fixed budget $K$ after validity and duplicate filtering. A stronger route-topology generator can replace this bank without changing the COWP certificate.

\section{Counterfactual Natural Alternatives}
\label{app:natural_alternatives}

\subsection{Observational Multimodal Branch}
\label{app:obs_branch}

The observational branch predicts a multimodal distribution conditioned on the full scene:
\begin{equation}
    p_\theta(\tau_i \mid \mathcal{S})
    =
    \sum_{m=1}^{M}
    \pi^m_i
    \mathcal{N}
    \left(
    \tau_i;
    \mu^m_i,
    \Sigma^m_i
    \right).
\end{equation}
We implement this branch as a transformer trajectory decoder over the shared burden-aware graph encoding. The branch is trained with a best-of-$M$ negative log-likelihood:
\begin{equation}
    \mathcal{L}_{\mathrm{obs}}
    =
    -\log
    \sum_{m=1}^{M}
    \pi^m_i
    \mathcal{N}
    \left(
    \tau^{\mathrm{gt}}_i;
    \mu^m_i,
    \Sigma^m_i
    \right).
\end{equation}
This branch provides data-distribution realism but is not sufficient alone because the ground-truth future may already include yielding to the ego.

\subsection{Ego-Neutral Counterfactual Branch}
\label{app:neutral_branch}

We define an ego-neutral intervention $\tau^{\mathrm{neu}}_e$ as an ego motion that removes pressure from the relevant conflict. Depending on the scene, it is generated by one of the following rules:
\begin{equation}
    \tau^{\mathrm{neu}}_e
    =
    \begin{cases}
        \tau^{\mathrm{yield}}_e, & \text{if ego approaches a merge or crossing conflict},\\
        \tau^{\mathrm{keep}}_e, & \text{if ego is not topologically involved},\\
        \tau^{\mathrm{stop}}_e, & \text{if ego faces stop/yield control}.
    \end{cases}
\end{equation}
The goal is not to predict what ego will do, but to remove the candidate coercive pressure.

We implement the neutral branch as a conditional denoising diffusion decoder. Let $\xi^0_i$ be a trajectory representation and $\xi^n_i$ its noisy version at diffusion step $n$. The denoising model predicts
\begin{equation}
    \epsilon_\theta
    =
    \epsilon_\theta
    \left(
    \xi^n_i,
    n,
    z_{\mathcal{G}},
    \eta(\tau^{\mathrm{neu}}_e)
    \right),
\end{equation}
where $z_{\mathcal{G}}$ is the graph encoding and $\eta(\tau^{\mathrm{neu}}_e)$ is an intervention embedding. The training loss is
\begin{equation}
    \mathcal{L}_{\mathrm{neutral}}
    =
    \mathbb{E}_{n,\epsilon}
    \left[
    \left\|
    \epsilon
    -
    \epsilon_\theta(\xi^n_i,n,z_{\mathcal{G}},\eta(\tau^{\mathrm{neu}}_e))
    \right\|_2^2
    \right].
\end{equation}
During training, the target can be either the observed future in non-conflict scenes or a pseudo-target generated by the priority-preserving branch in conflict scenes. This prevents the neutral predictor from learning that yielding to an aggressive ego is the default natural response.

\subsection{Priority-Preserving Rule Branch}
\label{app:priority_branch}

The priority-preserving branch constructs alternatives by solving a comfort-bounded motion primitive problem for the other agent:
\begin{equation}
\begin{aligned}
    \min_{\tau_i}\quad
    &
    J_{\mathrm{track}}(\tau_i)
    +
    J_{\mathrm{comfort}}(\tau_i)
    +
    J_{\mathrm{progress}}(\tau_i)\\
    \mathrm{s.t.}\quad
    &
    \tau_i(0)=x_i,\\
    &
    |a_i(t)|\le a^{\mathrm{comf}}_i,\quad
    |j_i(t)|\le j^{\mathrm{comf}}_i,\\
    &
    \tau_i(t)\in\mathcal{M}_{\mathrm{drivable}},\\
    &
    \mathrm{PriorityPreserved}(\tau_i,\mathcal{S})=1 .
\end{aligned}
\end{equation}
Priority preservation is defined by traffic controls, lane ownership, arrival order, and local topology. For example, a target-lane rear vehicle is allowed to maintain its gap and is not required to brake hard to accommodate an ego cut-in.

\subsection{Filtering and Calibration}
\label{app:natural_filtering}

The union in Eq.~\eqref{eq:natural_union} is filtered as follows:
\begin{equation}
    \mathcal{A}^{\mathrm{nat}}_i
    \leftarrow
    \left\{
    \tau_i
    \mid
    \mathrm{MapOK}(\tau_i)=1,\ 
    \mathrm{DynOK}(\tau_i)=1,\ 
    B_i(\tau_i;\mathcal{S},\tau^{\mathrm{neu}}_e)\le \beta_i(\mathcal{S}),\
    P_\theta(\tau_i)\ge p_{\min}
    \right\}.
\end{equation}
Option-space measure $\mu(\cdot)$ is implemented as the weighted sample mass:
\begin{equation}
    \mu(\mathcal{A})
    =
    \sum_{\tau_i^m\in\mathcal{A}} w_i^m,
    \qquad
    \sum_m w_i^m=1 .
\end{equation}
An optional extension would enlarge sampled alternatives using adaptive conformal radii,
\begin{equation}
    \mathcal{A}^{\mathrm{nat,cp}}_i
    =
    \bigcup_m
    \left\{
    \tau:
    d(\tau,\tau_i^m)\le q_{1-\delta}
    \right\},
\end{equation}
where $q_{1-\delta}$ is estimated on an independent calibration stream. This extension is not used in the reported implementation or experiments unless explicitly stated; no conformal coverage claim is made from the current natural-stage results.

\subsection{Natural Alternative Loss}
\label{app:natural_loss}

The natural alternative objective follows the source-restricted set, typed-prior, control, mode-usage, and root-identity decomposition in the main text. For every critical agent--root pair, the decoder first integrates the unconstrained control residual and computes the pre-projection multi-horizon ratio $\rho_{im}$ in Eq.~\eqref{eq:natural_mass_envelope}. With floor-smoothed detached mass $\widetilde p_{im}$, the exact optimized identity term is
\begin{equation}
    \mathcal{L}_{\mathrm{id}}
    =
    \operatorname{WM}_{i,m}
    \left(
      \left[
        \frac{\rho_{im}-\eta}{1-\eta}
      \right]_+^2;
      \operatorname{sg}(\widetilde p_{im})
    \right),
    \qquad 0<\eta<1,
\end{equation}
where $\operatorname{WM}$ is a normalized weighted mean over valid critical-agent roots. The loss is evaluated before emergency projection so that an over-envelope control retains a gradient toward the semantic interior. The wider emergency envelope is a separate numerical and physical safeguard: when activated, it scales the complete control sequence and re-integrates position, velocity, and heading consistently. Accordingly, component attribution evaluates both the exact squared-excess objective and probability mass outside the soft envelope; a change in the unoptimized global mean path ratio is reported only as a diagnostic.
The rule imitation term is
\begin{equation}
    \mathcal{L}_{\mathrm{prio}}
    =
    \min_m
    d(\tau_i^m,\tau_i^{\mathrm{prio}}),
\end{equation}
where $\tau_i^{\mathrm{prio}}$ is generated by Appendix~\ref{app:priority_branch}. The diversity term is
\begin{equation}
    \mathcal{L}_{\mathrm{div}}
    =
    -
    \frac{1}{M(M-1)}
    \sum_{m\neq n}
    \min
    \left(
    d(\tau_i^m,\tau_i^n),
    d_{\max}
    \right).
\end{equation}
This encourages coverage of multiple plausible non-coerced futures.

\section{Burden, Safe Responses, and Coercion Witnesses}
\label{app:burden_witness}

\subsection{Hybrid Burden Functional}
\label{app:burden}

The primitive burden functional combines rule-computable physical terms and a calibrated normative term:
\begin{equation}
\begin{aligned}
    B_i^{\mathrm{prim}}(\tau_i;\mathcal{S},\tau_e)
    =
    &
    w_a B^{\mathrm{acc}}_i
    +
    w_j B^{\mathrm{jerk}}_i
    +
    w_p B^{\mathrm{prog}}_i
    +
    w_r B^{\mathrm{risk}}_i \\
    &
    +
    w_n B^{\mathrm{norm}}_i .
\end{aligned}
\end{equation}
The acceleration burden is
\begin{equation}
    B^{\mathrm{acc}}_i
    =
    \max_t
    \left[
    \frac{|a_i(t)|-a^{\mathrm{comf}}_i}
    {a^{\mathrm{hard}}_i-a^{\mathrm{comf}}_i}
    \right]_+ .
\end{equation}
The jerk burden is
\begin{equation}
    B^{\mathrm{jerk}}_i
    =
    \frac{1}{T}
    \sum_t
    \left[
    \frac{|j_i(t)|-j^{\mathrm{comf}}_i}
    {j^{\mathrm{hard}}_i-j^{\mathrm{comf}}_i}
    \right]_+ .
\end{equation}
The progress burden compares the arrival time or longitudinal progress to a natural baseline:
\begin{equation}
    B^{\mathrm{prog}}_i
    =
    \frac{
    \left[
    T_i(\tau_i)-T_i(\tau_i^{\mathrm{nat}})
    \right]_+
    }{
    T_{\max}
    } .
\end{equation}
The risk exposure burden is
\begin{equation}
    B^{\mathrm{risk}}_i
    =
    \frac{1}{T}
    \sum_t
    \left(
    \left[
    \frac{\mathrm{TTC}_{\min}-\mathrm{TTC}_{ei}(t)}
    {\mathrm{TTC}_{\min}}
    \right]_+
    +
    \left[
    \frac{d_{\mathrm{safe}}(t)-d_{ei}(t)}
    {d_{\mathrm{safe}}(t)}
    \right]_+
    \right).
\end{equation}
Option preservation is not included inside $B_i^{\mathrm{prim}}$; it is enforced separately by Eq.~\eqref{eq:option_preservation}. This separation removes a circular definition in which the low-burden set depends on a burden term that itself depends on the low-burden set.

The normative burden is
\begin{equation}
    B^{\mathrm{norm}}_i
    =
    B^{\mathrm{rule}}_i
    +
    g_\theta(\mathcal{S},i,\tau_i,\tau_e),
\end{equation}
where $B^{\mathrm{rule}}_i$ encodes explicit priority violations and $g_\theta$ captures scene-conditioned conventions.

The burden threshold is context dependent but is not learned jointly with the certificate. In the evaluated implementation it is a frozen, pre-specified context table indexed by road-user type, local density/low-speed context, and priority relation:
\begin{equation}
    \beta_i(\mathcal{S})
    =\operatorname{clip}\!\left(\beta^{\mathrm{type}}_i+\Delta_{\mathrm{density}}(\mathcal{S})+\Delta_{\rho_i},\,\beta_{\min},\,\beta_{\max}\right).
\end{equation}
It is therefore independent of the certificate parameters and cannot be increased by the model to explain away coercive examples. A training-only empirical-quantile or conformal calibration is a reasonable future replacement, but the current implementation makes no finite-sample coverage claim for $\beta_i$.

\subsection{Ego-Conditioned Safe Response Set}
\label{app:safe_response}

For each ego candidate $\tau^k_e$, we generate agent response candidates by combining:
\begin{equation}
    \mathcal{R}_i(\tau^k_e)
    =
    \mathcal{R}^{\mathrm{pred}}_i(\tau^k_e)
    \cup
    \mathcal{R}^{\mathrm{opt}}_i(\tau^k_e)
    \cup
    \mathcal{R}^{\mathrm{emg}}_i(\tau^k_e).
\end{equation}
The predicted responses are sampled from an ego-conditioned decoder:
\begin{equation}
    \mathcal{R}^{\mathrm{pred}}_i(\tau^k_e)
    \sim
    p_\theta(\tau_i \mid \mathcal{S},\tau^k_e).
\end{equation}
The implementation does not solve an unrestricted continuous optimal-control problem for every candidate. Instead, Root-Conditioned Counterfactual Transport (RCOT) constructs an equal-budget timing-residual family $\mathcal{Q}_{im}(\tau_e^k)$ around each conflicting natural root and retains the lowest-burden safe member:
\begin{equation}
    b^\star_{ikm}
    =\min_{\tau_i\in\mathcal{Q}_{im}(\tau_e^k)}
    \left\{B_i^{\mathrm{prim}}(\tau_i;\mathcal{S},\tau_e^k):
    \mathrm{Collide}=0,\ \mathrm{DynOK}=1,\ \mathrm{MapOK}=1\right\},
\end{equation}
with the finite no-safe-response sentinel used only in implementation.  The family changes longitudinal timing---comfort deceleration, waiting, and mild recovery acceleration---while preserving the root polyline and topological maneuver.  Every conflicting root receives the same search budget before any compact response-slot truncation.  Emergency hard-braking or evasive-stop primitives are kept as a separate source and determine whether collision freedom exists only at high burden; their mere existence does not certify non-coercive feasibility.  A global top-$R$ response bank is retained only for auxiliary trajectory reconstruction and qualitative visualization.

The safe response set is
\begin{equation}
    \mathcal{R}^{\mathrm{safe}}_i(\tau^k_e)
    =
    \left\{
    \tau_i\in\mathcal{R}_i(\tau^k_e)
    \mid
    \mathrm{Unsafe}(\tau^k_e,\tau_i)=0
    \right\}.
\end{equation}
A response is low burden when $B_i^{\mathrm{prim}}\le\beta_i$; the retained mass follows Eq.~\eqref{eq:option_preservation}.

\subsection{Witness Pseudo-Label Construction}
\label{app:witness_supervision}

Because real datasets do not directly provide counterfactual labels, we construct pseudo-labels through counterfactual rollout.

For an ego candidate $\tau^k_e$ and agent $i$:
\begin{enumerate}
    \item Generate $\mathcal{A}^{\mathrm{nat}}_i$ using Appendix~\ref{app:natural_alternatives}.
    \item Roll out $\tau^k_e$ against each $\tau_i^{m}\in\mathcal{A}^{\mathrm{nat}}_i$.
    \item Compute the probability mass of ego-neutral-safe natural roots that become collision, near-miss, RSS-dangerous, or severe-margin violations under the ego candidate.
    \item Transport each conflicting root into its own response family and compute the minimum primitive burden among same-root safe responses.
    \item Assign a positive witness label only when conflicting mass exceeds $\delta_c$ and either retained root mass falls below $\alpha$ or conflict-conditioned tail burden exceeds $\gamma$.
\end{enumerate}

Formally,
\begin{equation}
    y^{\mathrm{wit}}_{ei,k}
    =
    \mathbf{1}\!\left[M^{\mathrm{conf}}_{ik}>\delta_c\right]
    \mathbf{1}\!\left[C_i(\tau^k_e)>\gamma\ \vee\ O_i(\tau^k_e)<\alpha\right].
\end{equation}
The mechanism token label is obtained by the largest normalized burden component:
\begin{equation}
    z^{\mathrm{gt}}_{ei,k}
    =
    \arg\max_{z\in\mathcal{Z}}
    B^z_i .
\end{equation}
Conflict intervals are obtained from the first and last time steps where swept volumes overlap, TTC falls below threshold, or safety distance is violated.

\subsection{Witness Loss}
\label{app:witness_loss}

The witness loss is
\begin{equation}
    \mathcal{L}_{\mathrm{exist}}
    =
    \mathrm{BCE}
    (
    \widehat{y}^{\mathrm{wit}}_{ei,k},
    y^{\mathrm{wit}}_{ei,k}
    ),
\end{equation}
\begin{equation}
    \mathcal{L}_{\mathrm{token}}
    =
    y^{\mathrm{wit}}_{ei,k}
    \cdot
    \mathrm{CE}
    (
    \widehat{z}_{ei,k},
    z^{\mathrm{gt}}_{ei,k}
    ),
\end{equation}
\begin{equation}
    \mathcal{L}_{\mathrm{burden}}
    =
    y^{\mathrm{wit}}_{ei,k}
    \cdot
    \left\|
    \widehat{B}_{ei,k}
    -
    B^{\mathrm{gt}}_{ei,k}
    \right\|_1,
\end{equation}
\begin{equation}
    \mathcal{L}_{\mathrm{conflict}}
    =
    y^{\mathrm{wit}}_{ei,k}
    \cdot
    \mathrm{SmoothL1}
    \left(
    [\widehat{t}_1,\widehat{t}_2],
    [t^{\mathrm{gt}}_1,t^{\mathrm{gt}}_2]
    \right).
\end{equation}

\subsection{Safe Response Supervision}
\label{app:response_loss}

For each sampled response $\tau_i^m$, we automatically label safety and low-burden status:
\begin{equation}
    y^{\mathrm{safe}}_{i,m}
    =
    \mathbf{1}
    [
    \mathrm{Unsafe}(\tau^k_e,\tau_i^m)=0
    ],
\end{equation}
\begin{equation}
    y^{\mathrm{low}}_{i,m}
    =
    \mathbf{1}
    [
    B_i(\tau_i^m;\mathcal{S},\tau^k_e)
    \le
    \beta_i(\mathcal{S})
    ].
\end{equation}
The response loss is
\begin{equation}
    \mathcal{L}_{\mathrm{resp}}
    =
    \mathrm{BCE}(\widehat{y}^{\mathrm{safe}},y^{\mathrm{safe}})
    +
    \mathrm{BCE}(\widehat{y}^{\mathrm{low}},y^{\mathrm{low}})
    +
    \lambda_B
    \|\widehat{B}-B\|_1 .
\end{equation}

\section{Plan Selection, Fallback, and Training}
\label{app:selection_training}

\subsection{Hard-First Inference}
\label{app:hard_first}

At inference, candidates are processed in three stages.

\paragraph{Stage 1: conventional safety.}
Reject candidates with predicted collision, road departure, dynamic infeasibility, or severe safety-margin violation:
\begin{equation}
    S_{\mathrm{coll}}(\tau^k_e)>\epsilon
    \Rightarrow
    \mathrm{Reject}(\tau^k_e).
\end{equation}

\paragraph{Stage 2: non-coercive feasibility.}
Let $\mathcal P^k$ denote the protected relations for candidate $k$.  Rule-observed \texttt{AgentPriority} and \texttt{EqualOrNegotiated} relations are hard members of $\mathcal P^k$ and cannot be removed by a learned priority score; learned relation probabilities are used only when the rule relation is unknown.  Reject candidates with a high-confidence protected witness:
\begin{equation}
    \max_{i\in\mathcal P^k}
    P(W^k_{ei}=1)
    >
    p_{\mathrm{hard}}
    \Rightarrow
    \mathrm{Reject}(\tau^k_e).
\end{equation}
The corresponding maximum over all critical agents is reported as a stricter diagnostic, not used as an unconditional hard veto.  For uncertain protected witnesses, we retain the candidate but add a penalty:
\begin{equation}
    J_{\mathrm{ncf}}(\tau^k_e)
    =
    \sum_{i\in\mathcal{I}^k}
    \left[
    \lambda_C [C_i(\tau^k_e)-\gamma]_+
    +
    \lambda_O [\alpha-O_i(\tau^k_e)]_+
    \right].
\end{equation}

\paragraph{Stage 3: ego utility.}
The final score is
\begin{equation}
    J(\tau^k_e)
    =
    J_{\mathrm{ego}}(\tau^k_e)
    +
    J_{\mathrm{ncf}}(\tau^k_e),
\end{equation}
where
\begin{equation}
    J_{\mathrm{ego}}
    =
    w_{\mathrm{prog}}J_{\mathrm{prog}}
    +
    w_{\mathrm{comfort}}J_{\mathrm{comfort}}
    +
    w_{\mathrm{route}}J_{\mathrm{route}}
    +
    w_{\mathrm{lane}}J_{\mathrm{lane}} .
\end{equation}
Candidates without high-confidence coercion are selected by minimum score.  Importantly, the semantic certificate set and the selector shortlist are different objects.  We denote
\begin{equation}
    \mathcal F_{\mathrm{cert}}
    =\{k:\text{$k$ satisfies conventional and protected RCOT constraints}\},
    \qquad
    \mathcal F_{\mathrm{short}}\subseteq\mathcal F_{\mathrm{cert}},
\end{equation}
where $\mathcal F_{\mathrm{short}}$ may be a bounded Pareto shortlist used only for efficient utility ranking.  Certificate acceptance, NCF recall, and coverage are computed from $\mathcal F_{\mathrm{cert}}$; they are never computed from the top-$K$ shortlist.

\subsection{Conservative Fallback}
\label{app:fallback}

If $\mathcal F_{\mathrm{cert}}$ is empty, the system enters an explicit \emph{uncertified fallback} state.  No primitive---including stopping or yielding---is assumed to be non-coercive, because abrupt deceleration can transfer burden to a close rear or merging vehicle.  Among conventionally valid candidates, fallback minimizes a robust composite risk consisting of a one-sided RCOT uncertainty bound, protected-rule risk, action/dynamics risk, pressure on nearby agents, and available rollout risk; ego utility is used only after these terms.  Controlled deceleration, yield-before-conflict, gap enlargement, creep, and emergency stop remain available primitives, but their names do not exempt them from the burden audit.

Every fallback event is recorded even when it returns a valid candidate index.  We report fallback frequency and fallback-conditioned PBTR separately.  The selected fallback is therefore a least-predicted-coercion emergency action, not a certified non-coercive plan, and it does not count toward certificate coverage.

\subsection{One-Sided Simultaneous Certificate Calibration (Planned Extension)}
\label{app:simultaneous_calibration}

A theorem-level extension calibrates the quantities used by the hard certificate rather than calibrating a single post-hoc candidate score.  For every protected candidate--agent--root tuple, let $\underline q_{ikm}$ be a lower confidence bound on same-root low-burden recovery and $\overline b^{\star}_{ikm}$ an upper confidence bound on minimum safe burden.  These induce a lower bound $\underline{\mathrm{OPR}}_i$ and an upper bound $\overline{\mathrm{CVaR}}_i$.  Candidate $k$ is accepted only if all protected relations satisfy
\begin{equation}
    \underline{\mathrm{OPR}}_i\ge \alpha_i,
    \qquad
    \overline{\mathrm{CVaR}}_i\le \gamma_i,
    \qquad \forall i\in\mathcal P^k.
\end{equation}
The calibration event must hold simultaneously over all protected relations inspected by the selector.  Under exchangeability this can be implemented with family-wise conformal or learn--then--test risk control; deployment under ego-induced interaction shift additionally requires weighted/robust calibration or an online risk monitor.  This subsection specifies the intended theoretical direction only.  The current experiments use learned uncertainty penalties and disjoint operating-point calibration and make no finite-sample simultaneous-coverage claim.

\subsection{Planning Loss}
\label{app:planning_loss}

The planning loss contains imitation, non-coercive ranking, and closed-loop outcome terms:
\begin{equation}
    \mathcal{L}_{\mathrm{plan}}
    =
    \lambda_{\mathrm{IL}}\mathcal{L}_{\mathrm{IL}}
    +
    \lambda_{\mathrm{rank}}\mathcal{L}_{\mathrm{rank}}
    +
    \lambda_{\mathrm{cl}}\mathcal{L}_{\mathrm{cl}}.
\end{equation}

The imitation loss selects the closest candidate to the expert trajectory:
\begin{equation}
    \mathcal{L}_{\mathrm{IL}}
    =
    \min_k
    d(\tau^k_e,\tau^{\mathrm{gt}}_e).
\end{equation}
The non-coercive ranking loss compares a non-coercive candidate with a coercive candidate:
\begin{equation}
    \mathcal{L}_{\mathrm{rank}}
    =
    \max
    \left(
    0,
    m
    +
    s_\theta(\tau^{\mathrm{coer}}_e)
    -
    s_\theta(\tau^{\mathrm{ncf}}_e)
    \right).
\end{equation}
The closed-loop loss is computed in simulation:
\begin{equation}
\begin{aligned}
    \mathcal{L}_{\mathrm{cl}}
    =
    &
    \lambda_{\mathrm{col}}\mathbf{1}[\mathrm{collision}]
    +
    \lambda_{\mathrm{near}}\mathbf{1}[\mathrm{near\text{-}miss}]
    +
    \lambda_{\mathrm{hb}}\mathbf{1}[\mathrm{other\ hard\ brake}]\\
    &
    +
    \lambda_{\mathrm{opt}} B^{\mathrm{option}}
    +
    \lambda_{\mathrm{prog}}J_{\mathrm{prog}}
    +
    \lambda_{\mathrm{comf}}J_{\mathrm{comfort}} .
\end{aligned}
\end{equation}
This loss prevents the model from relying solely on open-loop imitation, which may contain human or expert behaviors that are collision-free but coercive.

\subsection{Post-hoc Explanation}
\label{app:explanation}

For each rejected candidate, we report the highest-confidence coercion witness:
\begin{equation}
    \widehat{W}^{\star}
    =
    \arg\max_{\widehat{W}_{ei}}
    P(\widehat{W}_{ei}=1).
\end{equation}
The explanation includes the burdened agent, conflict region, conflict interval, mechanism token, burden magnitude, and option-preservation ratio. For example:
\begin{quote}
    Candidate \texttt{MergeAhead} is rejected because the rear vehicle in the target lane would need to perform \texttt{HB+GS}: hard braking and gap surrender. Its natural priority-preserving alternatives would enter a near-miss region under the ego candidate.
\end{quote}
This explanation is not merely post-hoc; the same token and burden quantities are also used for supervision and candidate rejection.

\printbibliography
\end{document}